%% ma: L11-B17-0074
%% lop: 11
%% chuong: 5
%% bai: 17
%% dang: 11.17.5
%% loai: TLN
%% muc: TH
%% dap_an: "-1"
%% nguon: "tai-lieu-hoc-tap-mon-toan-lop-11-theo-chuong-trinh-giao-duc-pho-thong-2018.pdf (Lớp Toán Thầy Duy, Phạm Lê Duy, THPT Chu Văn An), Chương 3 Bài 3, Phần 3, Câu 5 (tr.420)"
%% kien_thuc: "Lớp 11 Bài 17 (hàm số liên tục)"
%% trang_thai: cach-ly
%% ly_do: "Đáp án gốc là 3, nhưng với đề như in (x ≠ -2 và x = -2) đáp án là -1; đáp số 3 chỉ đúng nếu các điều kiện là x ≠ 2 và x = 2 (khi đó m = lim(x+1) = 3). Nghi lỗi in dấu; cần giáo viên xác nhận đề."
\begin{ex}
Tìm giá trị của $m$ để hàm số $f(x)=\begin{cases}\dfrac{x^2-x-2}{x-2}&\text{khi }x\ne-2\\m&\text{khi }x=-2\end{cases}$ liên tục trên tập xác định của chúng.
\shortans{-1}
\loigiai{\sloppy Với đề như in, $f(x)=x+1$ với $x\ne-2$, $x\ne2$ (hàm không xác định tại $x=2$), nên tập xác định là $\mathbb R\setminus\{2\}$. Liên tục tại $x=-2$ khi $m=\lim\limits_{x\to-2}(x+1)=-1$. Với $m=-1$, $f$ liên tục trên tập xác định.}
\end{ex}
